\chapter{Mouth Ulcers (Aphthous Ulcers)}

\section*{Definition}
Painful, shallow, round or oval sores on the oral mucosa, typically covered by a yellow-gray fibrinous membrane and surrounded by an erythematous halo, recurring at intervals (minor, major, or herpetiform types).

\section*{What Happens in Your Body}
Aphthous ulcers involve a localized immune-mediated inflammatory response targeting the oral epithelium. T-cell activation, tumor necrosis factor-alpha (TNF-alpha), and other immune signaling proteins cause epithelial breakdown. Trigger factors include minor trauma, stress, nutritional deficiencies, food sensitivities, and hormonal changes. Genetic predisposition is strong with positive family history in about 40\% of cases.

\section*{Causes}
\begin{itemize}
  \item Minor aphthous ulcers (most common, recur every few months)
  \item Major aphthous ulcers (larger, deeper, heal with scarring)
  \item Herpetiform ulcers (multiple clustered small ulcers)
  \item Trauma (biting cheek, sharp tooth, dental appliances, burns)
  \item Nutritional deficiencies (iron, folate, vitamin B12, zinc)
  \item Stress and anxiety (common trigger)
  \item Food sensitivities (chocolate, coffee, strawberries, nuts, cheese, tomatoes)
  \item Hormonal changes (premenstrual period)
  \item Autoimmune conditions (Crohn disease, Behcet disease, celiac disease)
  \item Immunosuppression (HIV, chemotherapy)
\end{itemize}

\section*{When to See a Doctor}
See your doctor if:
\begin{itemize}
  \item Symptoms are severe or getting worse over time
  \item Ulcers lasting more than 3 weeks, unusually large or painful ulcers, ulcers accompanied by systemic symptoms (fever, weight loss, rash, joint pain), or difficulty eating/drinking
  \item You have other concerning symptoms such as fever, weight loss, or bleeding
\end{itemize}

\section*{Related Symptoms}
\begin{itemize}
  \item Oral pain
  \item Difficulty swallowing
  \item Bad breath
  \item Swollen gums
  \item Sore throat
\end{itemize}
\textbf{Important:} If this symptom persists, your doctor will check for serious underlying causes.

\section*{Self-Care / Home Management}
\begin{itemize}
  \item Avoid known trigger foods (spicy, acidic, salty foods; chocolate, coffee, nuts, strawberries, tomatoes) during active outbreaks.
  \item Apply over-the-counter protective pastes (Orabase, triamcinolone acetonide in dental paste) directly to the ulcer to reduce pain and provide a protective barrier.
  \item Rinse with warm saline or alcohol-free chlorhexidine mouthwash 2-3 times daily to keep the area clean and reduce bacterial overgrowth.
  \item Use a soft-bristled toothbrush and avoid hard, sharp foods (toast, crisps) to prevent mechanical trauma that can trigger new ulcers.
\end{itemize}

\section*{How Your Doctor Will Evaluate You}
\begin{itemize}
  \item Visual inspection of ulcer number, size, location, and morphology (minor <1 cm, major >1 cm with scarring, herpetiform clusters).
  \item History of ulcer duration, frequency, triggers, and associated systemic symptoms (fever, rash, joint pain, gastrointestinal symptoms) to screen for underlying conditions.
  \item Blood tests for iron, ferritin, folate, vitamin B12, zinc, and inflammatory markers (CRP, ESR) when ulcers are frequent, severe, or atypical.
  \item Biopsy of the ulcer edge and base (with direct immunofluorescence) is indicated for persistent ulcers >3 weeks to exclude malignancy, vasculitis, or autoimmune disease.
\end{itemize}

\section*{Treatment Options}
\begin{itemize}
  \item Topical corticosteroids (hydrocortisone oromucosal tablets, clobetasol, or betamethasone mouthwash) reduce inflammation and pain; apply at the first sign of ulcer formation.
  \item For severe or frequent episodes, systemic therapies such as colchicine, dapsone, or low-dose thalidomide may be prescribed under specialist supervision.
  \item Addressing underlying nutritional deficiencies (iron, vitamin B12, folate, zinc) reduces recurrence in susceptible individuals.
  \item Chlorhexidine mouthwash or tetracycline mouthwash (not in children under 12) can reduce secondary bacterial infection and ulcer duration.
  \item Medications, lifestyle modifications, and physical therapies may be prescribed as appropriate.
  \item Follow-up ensures treatment effectiveness and allows timely adjustments to the care plan.
\end{itemize}

\section*{References}

\begin{thebibliography}{99}
\bibitem{harrison-aphthous} Longo DL, et al. \emph{Harrison's Principles of Internal Medicine}. 21st ed. McGraw-Hill; 2022. Chapter 45: Oral Manifestations of Disease.
\bibitem{uptodate-aphthous} Ship JA, et al. Recurrent aphthous stomatitis. In: UpToDate, Post TW (Ed), UpToDate, Waltham, MA; 2025.
\bibitem{chedid} Chedid V, et al. Aphthous stomatitis: a clinical review. \emph{JAMA}. 2023;329(15):1298-1308.
\bibitem{belenguer-guallar} Belenguer-Guallar I, et al. Treatment of recurrent aphthous stomatitis: a systematic review. \emph{J Oral Pathol Med}. 2024;53(2):101-114.
\bibitem{natah} Natah SS, et al. \emph{Burket's Oral Medicine}. 14th ed. Wiley-Blackwell; 2024. Chapter 5: Aphthous Ulcers.
\bibitem{rogers} Rogers RS, et al. Beh\c{c}et disease: an update on pathogenesis and management. \emph{Lancet Rheumatol}. 2023;5(7):e394-e407.
\end{thebibliography}
